\documentclass[12pt]{article}
\usepackage[utf8]{inputenc}
\usepackage[spanish]{babel}
\usepackage{amsmath, amssymb, amsthm}
\usepackage{hyperref}

\title{Practica 4\\Gramática para Analizador Sintáctico}
\author{Sarah Sophia Olivares García}
\date{}

\begin{document}

\maketitle

\section*{1. Conjuntos N, $\Sigma$, y símbolo inicial S}

Dada la gramática:
\[
\[
\begin{align*}
P = \{ & S \rightarrow \text{Expr} \ | \ \text{Asig}, \\
       & \text{Expr} \rightarrow \text{Term Expr}', \\
       & \text{Expr}' \rightarrow + \text{Term Expr}' \ | \ - \text{Term Expr}' \ | \ \varepsilon, \\
       & \text{Term} \rightarrow \text{Factor Term}', \\
       & \text{Term}' \rightarrow * \text{Factor Term}' \ | \ / \text{Factor Term}' \ | \ \varepsilon, \\
       & \text{Factor} \rightarrow \text{Num} \ | \ \text{Var} \ | \ (\text{Expr}) \ | \ -\text{Expr}, \\
       & \text{Num} \rightarrow \text{Entero Decimal}, \\
       & \text{Decimal} \rightarrow . \text{Entero} \ | \ \varepsilon, \\
       & \text{Entero} \rightarrow \text{Digito} \ | \ \text{Digito Entero}, \\
       & \text{Digito} \rightarrow 0 \ | \ 1 \ | \ldots \ | \ 9, \\
       & \text{Asig} \rightarrow \text{var Var = Expr}, \\
       & \text{Var} \rightarrow \text{Letra Pos}, \\
       & \text{Pos} \rightarrow \text{Var} \ | \ \varepsilon, \\
       & \text{Letra} \rightarrow \_ \ | \ a \ | \ b \ | \ldots \ | \ z \ | \ A \ | \ B \ | \ldots \ | \ Z
\}
\end{align*}
\]

\begin{itemize}
    \item $N$: \{S, Expr, Expr', Term, Term', Factor, Num, Decimal, Entero, Digito, Asig, Var, Pos, Letra\}
    \item $\Sigma$: \{+, -, *, /, (, ), ., var, =, \_, a-z, A-Z, 0-9\}
    \item Símbolo inicial: $S$
\end{itemize}

\section*{2. Eliminación de ambigüedad}

La gramática original presenta ambigüedad en la expresión ``Expr $\rightarrow$ Term Expr'``, ya que podría generar conflictos en la interpretación de operadores. Para eliminar la ambigüedad, utilizamos una jerarquía de operadores con precedencia y asociatividad definida:

\begin{itemize}
    \item Operadores de suma y resta tienen menor precedencia que los de multiplicación y división.
    \item La exponenciación es asociativa por la derecha.
\end{itemize}

Esto se puede manejar mediante reglas de precedencia en BYACC/J.

\section*{3. Eliminación de recursividad izquierda}

La regla $\text{Expr} \rightarrow \text{Term Expr}'$ no presenta recursividad izquierda directa o indirecta. Por tanto, no es necesario realizar modificaciones.

\section*{4. Factorización izquierda}

No se observa factor común entre las producciones. Por lo tanto, no se requiere factorización izquierda.

\section*{5. Nuevos conjuntos N y P}

Los conjuntos $N$ y $P$ permanecen iguales debido a que no hubo necesidad de eliminar recursividad o realizar factorizaciones.

\section*{6. Manejo de conflictos Shift/Reduce}

Los conflictos de shift/reduce se resolverán utilizando las declaraciones de precedencia y asociatividad en BYACC/J.

\section*{7. Definición en BYACC/J}

El archivo resultante en BYACC/J es el siguiente:

\begin{verbatim}
%{
  import java.lang.Math;
  import java.io.*;
%}

%token NUM VAR
%left '+' '-'
%left '*' '/'
%right '^'

%%

input: /* empty */
    | input line
;

line: '\n'
    | expr '\n' { System.out.println("Resultado: " + $1); }
;

expr: term
    | expr '+' term { $$ = $1 + $3; }
    | expr '-' term { $$ = $1 - $3; }
;

term: factor
    | term '*' factor { $$ = $1 * $3; }
    | term '/' factor { $$ = $1 / $3; }
;

factor: NUM
      | VAR
      | '(' expr ')' { $$ = $2; }
      | '-' factor %prec '-' { $$ = -$2; }
;

%%
\end{verbatim}

\section*{8. Ejemplo de entrada y salida}

Archivo de prueba (\texttt{prueba.txt}):
\begin{verbatim}
3 + 5
(2 * 6) - 4
\end{verbatim}

Salida esperada:
\begin{verbatim}
Resultado: 8
Resultado: 8
\end{verbatim}

\end{document}

